\section{势函数比较与直接应用}\label{app:section}

本节给出前述几何估计的直接应用。所有结论的方向都是“势函数或解控制分数次极大函数”；因此它们产生超水平集下界和资料集中限制，并不产生新的正则性上界。

\subsection{Riesz 势的方体比较}

对 $0<\beta<n$ 和 $f\ge0$，取未归一化 Riesz 势
\[
 I_\beta f(x)=\int_{\mathbb R^n}|x-y|^{\beta-n}f(y)\,dy.
\]

\begin{theorem}[Riesz 势比较及最佳常数]\label{app:riesz}
对每个 $f\ge0$，
\[
 I_\beta f(x)\ge n^{-(n-\beta)/2}M_{\beta,\mathcal Q}f(x),
\]
以及
\[
 I_\beta f(x)\ge\left(\frac{2}{\sqrt n}\right)^{n-\beta}
 M_{\beta,\mathcal Q_c}f(x).
\]
两个常数在“对所有非负 $f$、每一点成立”的意义下最佳。
\end{theorem}
\begin{proof}
若边长为 $\ell$ 的方体 $Q$ 包含 $x$，则 $|x-y|\le\sqrt n\ell$ 对 $y\in Q$ 成立。因 $\beta-n<0$，
\[
 \int_Q|x-y|^{\beta-n}f(y)\,dy
 \ge n^{-(n-\beta)/2}\ell^{\beta-n}\int_Qf.
\]
对允许的 $Q$ 取上确界即得第一个不等式。中心方体满足 $|x-y|\le\sqrt n\ell/2$，从而得到第二个不等式。

取 $x=0$，$y_0=(1,\ldots,1)$，以及
\[
 f_\varepsilon=\varepsilon^{-n}\mathbf1_{y_0+[-\varepsilon,0]^n}.
\]
支配收敛给出 $I_\beta f_\varepsilon(0)\to n^{-(n-\beta)/2}$。任何同时含 $0$ 且与该支集相交的非中心方体边长至少 $1-\varepsilon$，故 $M_{\beta,\mathcal Q}f_\varepsilon(0)\to1$。中心方体的最小边长趋于 $2$，故 $M_{\beta,\mathcal Q_c}f_\varepsilon(0)\to2^{\beta-n}$。这两个比值分别迫使所给常数不能增大。
\end{proof}

\begin{remark}\label{app:factor}
中心与非中心 Riesz 常数之比为 $2^{n-\beta}$。它与定理~\ref{mom:centered-weak} 中弱型常数右端的同一数值因子来自不同证明：前者来自核在中心方体上的直径，后者来自有限选择的重叠数。
\end{remark}

\subsection{Wolff 势的方体比较}

设 $1<p<n$，$\mu$ 为非负 Radon 测度，定义
\[
 W_{1,p}^{\mu}(x)=\int_0^\infty
 \left(\frac{\mu(B(x,r))}{r^{n-p}}\right)^{1/(p-1)}\frac{dr}{r}.
\]

\begin{theorem}[Wolff 势比较及最佳常数]\label{app:wolff}
令 $\gamma=(n-p)/(p-1)$。对每个 $x$，
\[
 W_{1,p}^{\mu}(x)\ge
 \frac{p-1}{n-p}\,n^{-\gamma/2}
 \bigl(M_{p,\mathcal Q}\mu(x)\bigr)^{1/(p-1)},
\]
并且
\[
 W_{1,p}^{\mu}(x)\ge
 2^\gamma\frac{p-1}{n-p}\,n^{-\gamma/2}
 \bigl(M_{p,\mathcal Q_c}\mu(x)\bigr)^{1/(p-1)},
\]
其中
\[
 M_{p,\mathcal B}\mu(x)=
 \sup_{Q\in\mathcal B(x)}|Q|^{p/n-1}\mu(Q).
\]
两个常数均最佳。
\end{theorem}
\begin{proof}
若 $Q\ni x$ 的边长为 $\ell$，令 $R=\sqrt n\ell$。在 Wolff 势中保留 $[R,\infty)$ 的积分，得到
\[
 W_{1,p}^{\mu}(x)\ge
 \mu(Q)^{1/(p-1)}\int_R^\infty r^{-\gamma}\frac{dr}{r}
 =\frac1\gamma n^{-\gamma/2}
 \bigl(\ell^{p-n}\mu(Q)\bigr)^{1/(p-1)}.
\]
取上确界即得第一个结论；中心情形中以 $R=\sqrt n\ell/2$ 代替即可。

若 $\mu=\delta_y$，则 $W_{1,p}^{\delta_y}(x)=\gamma^{-1}|x-y|^{-\gamma}$。选择 $x-y$ 的所有坐标绝对值相等时，$|x-y|=\sqrt n\|x-y\|_\infty$，而非中心和中心最小见证方体的边长分别为 $\|x-y\|_\infty$ 和 $2\|x-y\|_\infty$。直接代入表明两个常数均不能提高。
\end{proof}

\subsection{两个直接后果}

\begin{corollary}[分数 Laplace 资料的非退化区域]\label{app:fractional-laplace}
设 $0<\beta<\min\{2,n\}$，$E$ 是有界可测集，并令
\[
 u(x)=c_{n,\beta}\,I_\beta\mathbf1_E(x),\qquad
 c_{n,\beta}=\frac{\Gamma((n-\beta)/2)}{2^\beta\pi^{n/2}\Gamma(\beta/2)}.
\]
采用 Fourier 约定 $\widehat f(\xi)=\int_{\mathbb R^n}e^{-ix\cdot\xi}f(x)\,dx$，并以
\[
 \widehat{(-\Delta)^{\beta/2}f}(\xi)=|\xi|^\beta\widehat f(\xi)
\]
定义分数 Laplace 算子。此约定下，上述常数满足
$(-\Delta)^{\beta/2}u=\mathbf1_E$ 于缓增分布意义；该归一化及等价定义见 Kwa\'snicki \cite{Kwasnicki2017}。若 $E$ 是轴平行方体，则每个 $0<\alpha<1$ 满足
\[
 \bigl|\{u>c_{n,\beta}n^{-(n-\beta)/2}\alpha |E|^{\beta/n}\}\bigr|
 \ge (2\alpha^{-1/(n-\beta)}-1)^n|E|,
\]
并且
\[
 \bigl|\{u>c_{n,\beta}(2/\sqrt n)^{n-\beta}\alpha |E|^{\beta/n}\}\bigr|
 \ge \max\!\left\{(1-\alpha^{1/\beta})^n,
 (\alpha^{-1/(n-\beta)}-1)^n\right\}|E|.
\]
\end{corollary}
\begin{proof}
定理~\ref{app:riesz} 把相应方体光晕包含在 $u$ 的超水平集内。再使用非中心和中心方体的实际测试方体包含，即得两个不等式。
\end{proof}

\begin{corollary}[点态有界性排除资料集中]\label{app:concentration}
若 $f\ge0$ 且 $c_{n,\beta}I_\beta f(x)\le U<\infty$，则每个含 $x$ 的轴平行方体 $Q$ 满足
\[
 \int_Q f\le \frac{U}{c_{n,\beta}n^{-(n-\beta)/2}}\ell(Q)^{n-\beta}.
\]
对以 $x$ 为中心的方体，$n^{-(n-\beta)/2}$ 可改为较大的 $(2/\sqrt n)^{n-\beta}$。
\end{corollary}
\begin{proof}
由定理~\ref{app:riesz} 控制 $M_{\beta,\mathcal Q}f(x)$，再取单个方体作为该上确界的候选即可。
\end{proof}

\begin{corollary}[拟线性资料的集中必要条件]\label{app:plaplace}
设 $1<p<n$。假设非负 $p$-超调和函数 $u$ 已按 $\inf_{\mathbb R^n}u=0$ 归一化，Riesz 测度为 $\mu$，且在所考虑点满足可用的下势估计
\[
 u(x)\ge c_0W_{1,p}^{\mu}(x).
\]
若 $u(x)\le U$，则每个 $Q\ni x$ 有
\[
 \mu(Q)\le
 \left(\frac{U}{c_0\gamma^{-1}n^{-\gamma/2}}\right)^{p-1}
 \ell(Q)^{n-p}.
\]
中心方体中相应常数可提高 $2^\gamma$ 倍。
\end{corollary}
\begin{proof}
由定理~\ref{app:wolff} 解出 $M_{p,\mathcal Q}\mu(x)$，再取 $Q$ 为上确界的一个候选。这里的结论以所述下势估计和 Wolff 势在该点有限为前提；它并不对任意有限测度自动成立。
\end{proof}

\begin{remark}[应用边界]\label{app:boundary}
上述推论只给出了“解大于势下界”导出的必要条件和超水平集下界。Schauder、Calder\'on--Zygmund、梯度高可积性或 good--$\lambda$ 正则性仍需要反向分布控制、局部比较或容量估计；它们不能由这里的单向比较推出。有关分数极大分布函数和正则性接口的已有方法见 Nguyen--Tran \cite{NguyenTran2021}。
\end{remark}

\subsection{限制型范数、反向扩张与后续问题}

定理~\ref{gl:main}的全体中心光晕首项确定了限制型高矩的精确首项；进一步的统一估计还关系到近临界
$L^{r_p}\to L^p$ 示性资料常数。Bañuelos--Osękowski 的两权问题
\cite{BanuelosOsekowski2025}研究树状概率空间上固定指数的全函数范数；与本文可比较的
后续目标是欧氏示性资料的最优序列及其归一化形状。另一个方向由限制反向弱型
\cite{Lerner2026}提示：在非中心模型中识别对迭代封闭的近似基元素资料类，
证明相应的分数次反扩张与稳定性，再考察它如何进入局部分布估计。

本文已求得固定方体和任意固定锚定凸体的切向极限，并给出二维固定方体的高水平精确公式。
对全体可测输入，中心方体首项已经匹配；更精细的常数与一般几何基仍需研究。
二维半阶中，定理~\ref{pb:volume}将体积三阶剩余量夹在 $[2,10/3]$，
推论~\ref{mom:pair-bracket}将高矩次阶剩余量夹在 $[0,8]$；
定理~\ref{pb:strict-gap}已证明每个固定高阈值的两点上包络具有严格间隙；其端点衰减阶、极限存在性及上下界匹配仍未确定。具体问题包括：
\begin{enumerate}
\item 求高维非中心方体与强矩形的精确首项；确定中心方体的全体下一系数，并刻画可能并列的近极值构型。
\item 研究凸体基上对全体输入的匹配常数和仿射极值体。固定输入 $E=K$ 的联合切向体已由定理~\ref{cvx:full-tangent-limit}确定，对称半阶模型的椭球极值已由推论~\ref{half:ellipsoid-extremum}确定。
\item 在 $\beta\downarrow0$ 与 $\alpha\uparrow1$ 的耦合尺度中证明统一渐近；在二维精细公式的临界分数阶附近建立一致过渡。
\item 将近重合几何扩展到允许旋转、无界偏心率的几何族，或齐次群与 Ahlfors 正则空间。
\item 建立内禀一般测度模型的非中心与强矩形覆盖理论，研究度量球上一致弱型常数及极值结构。
\item 确定中心限制型高矩的下一系数及非中心精确首项，建立分数次反限制弱型迭代及其近极值理论。
\item 确定一般多输入几何平均模型的有限阈值常数，判断近极值资料是否渐近对角。
\item 构造带可控示性分解余项的反向 good--$\lambda$ 或容量估计，并在具体测度资料方程中检验定量正则性后果。
\end{enumerate}
